% Standalone problem document, generated from the adjacent README.md.
% Compile with XeLaTeX. Mathematical content is maintained in Markdown.
\documentclass[11pt,a4paper]{article}
\usepackage[fontsize=10.5pt]{scrextend}
\usepackage[margin=22mm,headheight=15pt,headsep=9mm,footskip=11mm]{geometry}
\usepackage{fontspec}
\setmainfont{texgyrepagella}[Extension=.otf,UprightFont=*-regular,BoldFont=*-bold,ItalicFont=*-italic,BoldItalicFont=*-bolditalic]
\setsansfont{texgyreheros}[Extension=.otf,UprightFont=*-regular,BoldFont=*-bold,ItalicFont=*-italic,BoldItalicFont=*-bolditalic]
\setmonofont{texgyrecursor}[Extension=.otf,UprightFont=*-regular,BoldFont=*-bold,ItalicFont=*-italic,BoldItalicFont=*-bolditalic,Scale=MatchLowercase]
\usepackage{amsmath,amssymb}
\usepackage{unicode-math}
\setmathfont{texgyrepagella-math.otf}
\usepackage{xcolor,microtype,enumitem,needspace,fancyhdr}
\usepackage{bookmark}
\definecolor{ink}{HTML}{17344B}
\definecolor{accent}{HTML}{197D80}
\definecolor{muted}{HTML}{596773}
\hypersetup{colorlinks=true,linkcolor=accent,urlcolor=accent,pdftitle={IE-15 - Exact
small-order growth factors for rook
pivoting},pdfauthor={Open Problems in Numerical Linear Algebra}}
\urlstyle{same}
\setlength{\parindent}{0pt}
\setlength{\parskip}{5pt plus 1pt minus 1pt}
\linespread{1.04}
\setlength{\emergencystretch}{3em}
\widowpenalty=10000
\clubpenalty=10000
\displaywidowpenalty=10000
\allowdisplaybreaks[1]
\setlist{nosep,leftmargin=1.5em}
\providecommand{\tightlist}{\setlength{\itemsep}{0pt}\setlength{\parskip}{0pt}}
\providecommand{\pandocbounded}[1]{#1}
\pagestyle{fancy}
\fancyhf{}
\fancyhead[L]{\sffamily\footnotesize\color{muted}OPEN PROBLEMS IN NUMERICAL LINEAR ALGEBRA}
\fancyhead[R]{\sffamily\bfseries\footnotesize\color{accent}IE-15}
\fancyfoot[L]{\parbox[b]{0.9\textwidth}{\sffamily\fontsize{8}{10}\selectfont\color{muted}Editorial ratings; a bounded search cannot certify that a problem remains unsolved.\\Literature check: 2026-09-11}}
\fancyfoot[R]{\sffamily\footnotesize\color{muted}\thepage}
\renewcommand{\headrulewidth}{0.3pt}
\renewcommand{\headrule}{\hbox to\headwidth{\color{accent}\leaders\hrule height \headrulewidth\hfill}}
\makeatletter
\renewcommand{\section}{\Needspace{5\baselineskip}\@startsection{section}{1}{0pt}{12pt plus 2pt minus 2pt}{4pt}{\sffamily\bfseries\large\color{ink}}}
\renewcommand{\subsection}{\Needspace{4\baselineskip}\@startsection{subsection}{2}{0pt}{9pt plus 1pt minus 1pt}{3pt}{\sffamily\bfseries\normalsize\color{ink}}}
\makeatother
\setcounter{secnumdepth}{0}
\begin{document}
{\sffamily\small\color{muted}Linear systems and elimination\par}
\vspace{3pt}
{\raggedright\hyphenpenalty=10000\exhyphenpenalty=10000\sffamily\bfseries\fontsize{21}{24}\selectfont\color{ink}Exact
small-order growth factors for rook pivoting\par}
\vspace{7pt}
{\sffamily\small\textbf{Difficulty:} hard\quad\textbf{Importance:} interesting
to specialist\par}
{\sffamily\small\bfseries\color{accent}Status: Open\par}
\vspace{4pt}
{\color{accent}\hrule height 0.8pt}
\vspace{6pt}
\subsection{Related order-five bound -
2026-09-11}\label{related-order-five-bound---2026-09-11}

Matthew J. Colbrook submitted a recovered rational \(5\times5\) matrix
with an admissible rook path of growth \(893/131\). The
\href{https://github.com/ajt60gaibb/OpenProblemsInNLA/blob/main/references/colbrook-recovered-2026-09-11/submitted/research/rook_partial.md}{complete
construction},
\href{https://github.com/ajt60gaibb/OpenProblemsInNLA/blob/main/references/colbrook-recovered-2026-09-11/verification/fresh-rook-results.json}{exact
rerun}, and
\href{https://github.com/ajt60gaibb/OpenProblemsInNLA/blob/main/references/colbrook-recovered-2026-09-11/verification/reviews/rook-review.md}{independent
review} verify this finite lower bound, with ties allowed. It concerns
order five and determines neither requested order-three nor order-four
constant. \textbf{IE-15 remains Open}, with its target and ratings
unchanged. See the
\href{https://github.com/ajt60gaibb/OpenProblemsInNLA/blob/main/references/colbrook-recovered-2026-09-11/README.md}{submission
record} for the Cambridge affiliation, AI-assistance disclosure and
verification limits.

\textbf{Rating rationale:} Hard reflects two focused finite-dimensional
extremal constants within an established pivoting model; specialist
impact concerns exact small-order rook-pivoting behavior.

\subsection{Problem statement}\label{problem-statement}

All arithmetic is exact and matrices are real and nonsingular. At each
step of Gaussian elimination with rook pivoting, select a nonzero entry
maximal in absolute value both in its row and in its column of the
active matrix. Move it to the active \((1,1)\) position by row and
column interchanges, and form the trailing Schur complement.

For an admissible path \(\pi\) on \(A\in\mathbb R^{n\times n}\), let
\(S_1=A,\ldots,S_n\) be the active matrices and define

\[
\rho(A,\pi)=\frac{\max_{1\le k\le n}\|S_k\|_{\max}}{\|A\|_{\max}},
\qquad \|M\|_{\max}=\max_{i,j}|m_{ij}|,
\] \[
g_{\mathrm{RP}}(n)=\sup_{\substack{A\in\mathbb R^{n\times n}\\\det A\ne0}}\
\sup_{\pi\text{ permitted by rook pivoting}}\rho(A,\pi).
\]

Determine the two exact constants \(g_{\mathrm{RP}}(3)\) and
\(g_{\mathrm{RP}}(4)\), with matching upper and lower bounds. All
admissible rook choices and ties are included. The source singles out
dimensions at most four; \(g_{\mathrm{RP}}(1)=1\) and
\(g_{\mathrm{RP}}(2)=2\) are known. The two unknown constants are
counted as one problem.

Rook pivoting balances pivot-search cost and growth control. Sharp
small-order constants would sharpen its finite-dimensional stability
theory. General asymptotic growth estimates do not determine them; IE-11
concerns a different pivoting rule.

\newpage
\subsection{References}\label{references}

N. J. Higham,
\href{https://doi.org/10.1137/1.9780898718027}{\emph{Accuracy and
Stability of Numerical Algorithms}, second edition}, SIAM (2002),
Problem 9.18, p.~193, and definition (9.15), p.~169. A. Edelman and J.
Urschel, \href{https://doi.org/10.1137/23M1571903}{\emph{Some New
Results on the Maximum Growth Factor in Gaussian Elimination}}, SIMAX 45
(2024), \href{https://arxiv.org/html/2303.04892v4}{§6}. R. Shah and J.
Urschel, \href{https://arxiv.org/html/2608.19189v4}{\emph{Entry growth
in Gaussian elimination}}, August 31, 2026 revision, Theorem 1.6 (=6.2).

\subsection{Earlier status check ---
2026-09-08}\label{earlier-status-check-2026-09-08}

The book's question and real-field supremum were checked directly.
Searches for \textit{rook pivoting 3x3}, \textit{rook pivoting 2.9},
\textit{rook pivoting 4.16}, and
\textit{rook pivoting maximum growth exact value} found no exact
small-order solution. The cited 2024 large-order construction and 2026
asymptotic result do not settle this pair. The book's decimals 2.9 and
4.16 are historical numerical search outcomes, not exact values. No
recent source explicitly reaffirming this particular question's openness
was located.

\subsection{Audit update --- 2026-09-10}\label{audit-update-2026-09-10}

Rechecked Higham's
\href{https://pages.stat.wisc.edu/~bwu62/771/hingham2002.pdf}{Problem
9.18}, p.~193. Searches for exact order-three/order-four rook growth and
the \href{https://arxiv.org/html/2608.19189v4}{2026 general pivoting
bounds} found no exact values. Difficulty is reduced to hard to reflect
the focused finite-order target, although the global upper-bound proofs
may still be substantial. The open verdict remains bounded by a
historical explicit source.
\end{document}
